\chapter{Покажчик пунктів реєстрів}\label{app:index}

Тут зібрано \emph{кожен} пункт усіх реєстрів проєкту: батьківських
(\texttt{Our try/00-decisions}: E, R, C, Q, D \cite{RegEst,RegConj,RegDec,RegQ}), журналу
верифікації (V-\,\cite{RegVerif}), реєстру візуалізації (VE, VR, VC, VQ \cite{RegViz}) і трьох
пунктів бекенда котушок (U1--U3), які в \texttt{tokamak-3d-viz/docs/COIL-BACKEND.md} не мають
номерів (номери U дав путівник; див.\ розд.~\ref{ch:g09}).

\paragraph{Як зроблено таблицю.} Тіло таблиці генерує скрипт
\texttt{guide/tools/build\_index.py}. Він \emph{парсить} реєстри, тож ідентифікатори, статуси
й номери рядків беруться з самих файлів (стан на 2026-10-07), а не переписуються вручну. Вручну
задано лише три стовпці: суть (до 12 слів), посилання на код і розділ путівника. Скрипт
завершується з помилкою, якщо в реєстрі з'явився пункт без рядка в покажчику (або навпаки), а
також якщо посилання на код указує на файл чи функцію, яких немає. Після зміни реєстрів:
\verb|cd guide && python3 tools/build_index.py|.

{\raggedright\paragraph{Стовпці.} «Де в реєстрі»~--- файл і номер рядка (р.), з якого починається запис. «Код»~---
\texttt{файл:функція} з тими самими префіксами, що й у фрагментах коду путівника
(\texttt{ourtry/}~= \texttt{Our try/}; \texttt{viz/}~= \texttt{tokamak-3d-viz/};
\texttt{gs/}~= \texttt{Tokamak-GS-solver/}; \texttt{fec/}~= \texttt{fusion equilibrium
challenge/}; \texttt{manualcode/}~= \texttt{manual/code/}); «---» означає, що пункт прочитано в джерелі або
він організаційний. «Розділ»~--- де пункт розповідається; «лише тут» позначає технічні пункти
(рендер Blender, перевірки чисел Sophelio тощо), які путівник свідомо не розбирає.
Стрілка $\to$ у стовпці «суть» вказує на пов'язаний пункт іншого реєстру.\par}

\makeatletter\newcommand\gIndexRows{\@@input appendix/index_rows.tex }\makeatother
{\footnotesize
\setlength{\tabcolsep}{2.5pt}
\renewcommand{\arraystretch}{1.05}
\begin{longtable}{@{}>{\raggedright\arraybackslash}p{1.05cm}
                    >{\raggedright\arraybackslash}p{4.75cm}
                    >{\raggedright\arraybackslash}p{2.35cm}
                    >{\raggedright\arraybackslash}p{2.55cm}
                    >{\raggedright\arraybackslash}p{3.4cm}
                    >{\raggedright\arraybackslash}p{1.35cm}@{}}
\toprule
\textbf{ID} & \textbf{Суть} & \textbf{Статус} & \textbf{Де в реєстрі} & \textbf{Код} & \textbf{Розділ}\\
\midrule
\endfirsthead
\toprule
\textbf{ID} & \textbf{Суть} & \textbf{Статус} & \textbf{Де в реєстрі} & \textbf{Код} & \textbf{Розділ}\\
\midrule
\endhead
\midrule
\multicolumn{6}{r}{\emph{продовження на наступній сторінці}}\\
\endfoot
\bottomrule
\endlastfoot
\gIndexRows
\end{longtable}
}

\paragraph{Підсумок за реєстрами.} E1--E40~--- 40; R1--R13~--- 13; C1--C10~--- 10 (C8 закрито~→ R9, C1~→ R10, C2~→ R11, C5~→ R13; половина C4~→ R12);
Q1--Q13~--- 13 (Q11, Q12 закриті); D1--D9~--- 9; журнал верифікації V-A1\ldots V-C4~--- 17;
VE1--VE33 разом із VE11a~--- 34; VR1--VR16~--- 16; VC1--VC7~--- 7; VQ1--VQ5~--- 5; U1--U3~--- 3.
Разом \textbf{167} пунктів (лічильники друкує той самий скрипт). Шість відкинутих варіантів і
три визнані ризики з \texttt{DECISIONS.md} номерів не мають і в покажчик не входять; їх
розібрано в розд.~\ref{ch:g01}.
